\chapter{Nhà Nghèo Không Nghèo Ý Chí}
\markboth{Nhà Nghèo Không Nghèo Ý Chí}{A Poor Home Is Not a Poor Will}

\section{Áo Cũ Nhưng Bài Vẫn Đủ}
\begin{truyen}{Áo Cũ Nhưng Bài Vẫn Đủ}{Old Clothes but Homework Still Done}
\chuhoa{C}{ó em đi học với đôi dép cũ, áo hơi sờn, cặp không mới.}
\textit{(Some students come to school in worn sandals, faded shirts, and old bags.)}

Nhưng bài em làm rất ngay ngắn, tập em giữ rất kỹ.
\textit{(Yet their work is neat and their notebooks are carefully kept.)}

Nhìn một đứa trẻ như vậy mới thấy nghèo không đồng nghĩa với buông.
\textit{(Seeing such a child reminds us that poverty does not automatically mean giving up.)}

Có những gia đình ít tiền nhưng dạy con giữ lòng tự trọng rất rõ.
\textit{(Some families have little money yet teach dignity very clearly.)}

\end{truyen}

\begin{baihoc}
Thiếu điều kiện không có nghĩa là thiếu phẩm chất.
\textit{(Lack of resources does not mean lack of character.)}
\end{baihoc}

\begin{ghinhoanh}
\textbf{Short English to remember:}

Lack of resources does not mean lack of character.

\end{ghinhoanh}

\begin{tuhocnhanh}
\textbf{Gợi ý để dạy và tự nghĩ / Teaching and reflection prompts}

1. Điều gì đẹp trong hình ảnh này? \textit{What is beautiful in this picture?}

2. Vì sao không nên nhìn học sinh bằng đồ dùng bên ngoài? \textit{Why should students not be judged by outward things?}

3. Mình có đang đủ tinh tế với những em như vậy không? \textit{Are we sensitive enough to students like this?}
\end{tuhocnhanh}
\ngancach

\section{Nghỉ Học Một Buổi Vì Tiền}
\begin{truyen}{Nghỉ Học Một Buổi Vì Tiền}{Missing School Because of Money}
\chuhoa{C}{ó em nghỉ học không phải vì lười, mà vì hôm đó nhà không xoay kịp một khoản rất nhỏ.}
\textit{(Some students miss school not because of laziness, but because home cannot manage a very small expense that day.)}

Với người khác, số tiền ấy không lớn.
\textit{(To others, that amount seems small.)}

Nhưng với một nhà đang chật vật, cái nhỏ đó đủ làm một đứa trẻ ngồi lại ở nhà.
\textit{(But in a struggling home, that small amount is enough to keep a child away from class.)}

Có những khoảng cách trong giáo dục được tạo ra bởi những thứ rất bé mà rất đau.
\textit{(Some educational gaps are created by things that look very small but hurt deeply.)}

\end{truyen}

\begin{baihoc}
Nhiều rào cản của học sinh không nằm ở ý chí, mà nằm ở hoàn cảnh.
\textit{(Many student barriers are not about willpower, but about circumstance.)}
\end{baihoc}

\begin{ghinhoanh}
\textbf{Short English to remember:}

Many student barriers come from circumstance, not from willpower.

\end{ghinhoanh}

\begin{tuhocnhanh}
\textbf{Gợi ý để dạy và tự nghĩ / Teaching and reflection prompts}

1. Điều gì đang cản trở em này? \textit{What is holding this student back?}

2. Vì sao người dạy cần biết hoàn cảnh thật? \textit{Why should teachers understand real circumstances?}

3. Một hỗ trợ nhỏ có thể đổi được gì? \textit{What can a small support change?}
\end{tuhocnhanh}
\ngancach

\section{Xin Nộp Chậm Không Phải Vì Coi Thường}
\begin{truyen}{Xin Nộp Chậm Không Phải Vì Coi Thường}{Asking for More Time Is Not Always Disrespect}
\chuhoa{N}{hiều em ngại nói, nên khi xin nộp bài chậm thường bị hiểu ngay là lười hoặc vô trách nhiệm.}
\textit{(Many students struggle to explain, so asking for more time is quickly read as laziness or irresponsibility.)}

Nhưng có em buổi tối phải phụ việc, trông em nhỏ, hoặc sống trong một nhà không có lấy góc yên để học.
\textit{(Yet some spend evenings working, caring for siblings, or living in homes without a quiet corner to study.)}

Không phải lời xin thêm thời gian nào cũng đáng tin, nhưng cũng không phải lời nào cũng đáng nghi.
\textit{(Not every request for time is trustworthy, but not every one deserves suspicion.)}

Người dạy tốt thường vừa giữ nguyên tắc vừa còn chỗ cho hoàn cảnh thật.
\textit{(A good teacher keeps standards while still leaving room for real circumstances.)}

\end{truyen}

\begin{baihoc}
Kỷ luật và lòng hiểu người không nhất thiết phải đứng ở hai phía đối nhau.
\textit{(Discipline and understanding do not always have to stand on opposite sides.)}
\end{baihoc}

\begin{ghinhoanh}
\textbf{Short English to remember:}

Discipline and understanding do not always stand on opposite sides.

\end{ghinhoanh}

\begin{tuhocnhanh}
\textbf{Gợi ý để dạy và tự nghĩ / Teaching and reflection prompts}

1. Người dạy cần giữ điều gì và mềm ở điều gì? \textit{What should a teacher hold firm, and where should the teacher soften?}

2. Vì sao không nên kết luận quá nhanh? \textit{Why should we avoid quick conclusions?}

3. Em học sinh cần học thêm điều gì về trách nhiệm? \textit{What does the student still need to learn about responsibility?}
\end{tuhocnhanh}
\ngancach

\section{Tự Trọng Trong Cái Nghèo}
\begin{truyen}{Tự Trọng Trong Cái Nghèo}{Dignity Inside Poverty}
\chuhoa{C}{ó em nghèo thật nhưng rất ít khi xin xỏ quá lời. Các em nhận sự giúp đỡ với một sự dè dặt rất rõ.}
\textit{(Some students are truly poor but rarely ask too much. They receive help with visible hesitation.)}

Các em không muốn bị thương hại, dù đang cần giúp.
\textit{(They do not want pity, even when they need support.)}

Cái tự trọng đó vừa làm người lớn thương, vừa nhắc người lớn phải giúp cho đúng cách.
\textit{(That dignity both moves adults and reminds them to help in the right way.)}

Giúp một đứa trẻ không chỉ là đưa cho nó cái thiếu. Còn là giữ cho nó không thấy mình bị hạ xuống.
\textit{(Helping a child is not only giving what is missing. It is also preserving the child’s sense of worth.)}

\end{truyen}

\begin{baihoc}
Sự giúp đỡ đẹp là sự giúp đỡ không làm người nhận thấy mình bé đi.
\textit{(Beautiful help does not make the receiver feel smaller.)}
\end{baihoc}

\begin{ghinhoanh}
\textbf{Short English to remember:}

Beautiful help does not make the receiver feel smaller.

\end{ghinhoanh}

\begin{tuhocnhanh}
\textbf{Gợi ý để dạy và tự nghĩ / Teaching and reflection prompts}

1. Vì sao lòng tự trọng này đáng được giữ? \textit{Why is this dignity worth protecting?}

2. Giúp đúng cách khác giúp cho xong ở đâu? \textit{How is right help different from convenient help?}

3. Mình có đang giúp mà vô tình làm đau không? \textit{Could our help be hurting unintentionally?}
\end{tuhocnhanh}
\ngancach

\section{Đi Học Với Bụng Chưa No}
\begin{truyen}{Đi Học Với Bụng Chưa No}{Coming to School on an Empty Stomach}
\chuhoa{C}{ó em vào tiết một mà mắt đã đờ ra. Hỏi nhỏ mới biết sáng nay chưa ăn gì.}
\textit{(Some students enter first period already looking drained. A quiet question reveals they have eaten nothing.)}

Một cái bụng rỗng làm chữ nghĩa khó vào hơn rất nhiều.
\textit{(An empty stomach makes learning much harder to receive.)}

Nhiều người lớn nói về ý chí, nhưng có những thứ cơ bản đến mức thiếu nó thì ý chí cũng bị kéo xuống.
\textit{(Adults talk about willpower, but some conditions are so basic that without them even willpower gets pulled down.)}

Có những bài học bắt đầu từ một ổ bánh mì.
\textit{(Some lessons begin with a piece of bread.)}

\end{truyen}

\begin{baihoc}
Muốn nói về học tập, đôi khi phải nhìn lại cả những điều rất căn bản của đời sống.
\textit{(If we want to talk about learning, we sometimes must first face the most basic parts of living.)}
\end{baihoc}

\begin{ghinhoanh}
\textbf{Short English to remember:}

Some lessons begin with basic care.

\end{ghinhoanh}

\begin{tuhocnhanh}
\textbf{Gợi ý để dạy và tự nghĩ / Teaching and reflection prompts}

1. Điều cơ bản nào đang thiếu ở đây? \textit{What basic need is missing here?}

2. Vì sao không thể chỉ nói về ý chí? \textit{Why can’t we talk only about willpower?}

3. Một giúp đỡ nhỏ lúc này có thể là gì? \textit{What small help could matter right now?}
\end{tuhocnhanh}
\ngancach

\section{Đi Bộ Xa Nhưng Vẫn Đúng Giờ}
\begin{truyen}{Đi Bộ Xa Nhưng Vẫn Đúng Giờ}{Walking Far Yet Still Arriving on Time}
\chuhoa{C}{ó em nhà xa, phương tiện không thuận, trời mưa vẫn đi học gần như không nghỉ.}
\textit{(Some students live far away with difficult transport, yet still come to school almost every day.)}

Điều đó nhìn quen rồi nên người ta dễ coi là bình thường.
\textit{(Because it happens quietly, people begin to treat it as normal.)}

Nhưng nhiều nỗ lực lớn của học sinh lại nằm ở những việc rất ít được kể công như vậy.
\textit{(Yet many of a student’s biggest efforts live in these ordinary, rarely praised routines.)}

Có em đã cố rất nhiều trước cả khi ngồi xuống bàn học.
\textit{(Some children have already tried very hard before they even sit down to learn.)}

\end{truyen}

\begin{baihoc}
Đừng chỉ nhìn lúc học sinh ngồi trong lớp. Hãy nhớ có em đã vượt một đoạn đường dài để tới được đó.
\textit{(Do not only look at students once they sit in class. Some have already crossed a long road just to get there.)}
\end{baihoc}

\begin{ghinhoanh}
\textbf{Short English to remember:}

Some students have already crossed a long road before class even begins.

\end{ghinhoanh}

\begin{tuhocnhanh}
\textbf{Gợi ý để dạy và tự nghĩ / Teaching and reflection prompts}

1. Nỗ lực nào đang bị coi là bình thường quá mức? \textit{What effort is being treated as too ordinary?}

2. Vì sao cần nhìn cả quãng đường tới lớp? \textit{Why should we see the road to school too?}

3. Mình có đang biết ơn đủ những nỗ lực thầm lặng không? \textit{Are we grateful enough for quiet effort?}
\end{tuhocnhanh}
\ngancach

\section{Không Muốn Bạn Biết Nhà Mình Nghèo}
\begin{truyen}{Không Muốn Bạn Biết Nhà Mình Nghèo}{Not Wanting Classmates to Know the Family Is Poor}
\chuhoa{C}{ó em rất ngại các hoạt động đụng tới chuyện tiền, chuyện nhà, hoặc chuyện mời bạn về chơi.}
\textit{(Some students become uneasy whenever class activities touch money, home life, or inviting friends over.)}

Các em sợ bị nhìn khác đi.
\textit{(They fear being seen differently.)}

Nghèo không chỉ làm trẻ thiếu thứ này thứ kia. Nó còn làm trẻ phải giấu.
\textit{(Poverty does not only create material lack. It also teaches children to hide.)}

Có những em học cách cười bình thường để không ai chạm vào phần các em đang xấu hổ.
\textit{(Some students learn to smile normally so no one touches the part they feel ashamed of.)}

\end{truyen}

\begin{baihoc}
Sự tinh tế của lớp học có thể giữ được lòng tự trọng cho một học sinh nghèo.
\textit{(Classroom sensitivity can protect the dignity of a poor student.)}
\end{baihoc}

\begin{ghinhoanh}
\textbf{Short English to remember:}

Classroom sensitivity can protect a poor student’s dignity.

\end{ghinhoanh}

\begin{tuhocnhanh}
\textbf{Gợi ý để dạy và tự nghĩ / Teaching and reflection prompts}

1. Học sinh này đang giấu điều gì? \textit{What is this student trying to hide?}

2. Vì sao sự tinh tế của lớp học quan trọng? \textit{Why is classroom sensitivity important here?}

3. Người dạy có thể tránh những cách tổ chức nào? \textit{What kinds of activities should the teacher handle more carefully?}
\end{tuhocnhanh}
\ngancach

\section{Ý Chí Không Ồn Ào}
\begin{truyen}{Ý Chí Không Ồn Ào}{Willpower Is Often Quiet}
\chuhoa{N}{hiều em không kể khổ, không nói cố gắng nhiều, không xin được thương. Các em chỉ đi học đều, làm bài đều, và chịu khó đều.}
\textit{(Many students do not talk about hardship, do not advertise effort, and do not ask for pity. They simply keep showing up, keep working, and keep enduring.)}

Ý chí của các em không ồn ào nên cũng ít được khen.
\textit{(Their willpower is quiet, so it is praised less often.)}

Nhưng chính kiểu bền này mới làm người ta nể lâu.
\textit{(Yet this kind of quiet endurance is what earns lasting respect.)}

Có những đứa trẻ nghèo không làm mình thương bằng lời. Chúng làm mình nể bằng cách sống.
\textit{(Some poor children move us not by words, but by the way they live.)}

\end{truyen}

\begin{baihoc}
Ý chí đẹp nhất nhiều khi không phải ý chí nói ra, mà là ý chí được lặp lại mỗi ngày.
\textit{(The most beautiful willpower is often not spoken. It is repeated daily.)}
\end{baihoc}

\begin{ghinhoanh}
\textbf{Short English to remember:}

The most beautiful willpower is often repeated daily, not loudly declared.

\end{ghinhoanh}

\begin{tuhocnhanh}
\textbf{Gợi ý để dạy và tự nghĩ / Teaching and reflection prompts}

1. Ý chí ở đây biểu hiện bằng điều gì? \textit{How does willpower show itself here?}

2. Vì sao ý chí thầm lặng dễ bị quên? \textit{Why is quiet endurance easily forgotten?}

3. Mình có đang khen đúng kiểu nỗ lực này không? \textit{Are we praising this kind of effort enough?}
\end{tuhocnhanh}
